% Part: propositional-logic
% Chapter: propositional-logic
% Section: preliminaries

\documentclass[../../../include/open-logic-section]{subfiles}

\begin{document}

\olfileid{pl}{syn}{pre}

\olsection{مُڈھلیاں گلاں}

\begin{thm}[\emph{!!{formula}s اُتے استقرا دا اصول}]
  \ollabel{thm:induction}
  جے کوئی خاصیت~$P$ سارے ایٹمی !!{formula}s لئی پوری ہُندی اے تے
  اینج اے کہ
  \begin{enumerate}
    \tagitem{prvNot}{ایہ $\lnot !A$ لئی پوری ہُندی اے جدوں وی ایہ
      ~$!A$ لئی پوری ہُندی اے؛}{}
    \tagitem{prvAnd}{ایہ $(!A \land !B)$ لئی پوری ہُندی اے
      جدوں وی ایہ $!A$ تے~$!B$ لئی پوری ہُندی اے؛}{}
    \tagitem{prvOr}{ایہ $(!A \lor !B)$ لئی پوری ہُندی اے
      جدوں وی ایہ $!A$ تے~$!B$ لئی پوری ہُندی اے؛}{}
    \tagitem{prvIf}{ایہ $(!A \lif !B)$ لئی پوری ہُندی اے
      جدوں وی ایہ $!A$ تے~$!B$ لئی پوری ہُندی اے؛}{}
    \tagitem{prvIff}{ایہ $(!A \liff !B)$ لئی پوری ہُندی اے
      جدوں وی ایہ $!A$ تے~$!B$ لئی پوری ہُندی اے؛}{}
  \end{enumerate}
  تاں $P$ سارے !!{formula}s لئی پوری ہُندی اے۔
\end{thm}

\begin{proof}
  $S$ نوں اوہناں سارے !!{formula}s دی جماعت منّو جنہاں وچ
  خاصیت~$P$ اے۔ صاف اے کہ $S \subseteq \Frm[L_0]$۔ $S$،
  \olref[fml]{defn:formulas} دیاں ساریاں شرطاں پوری کردا اے: ایہدے وچ
  سارے ایٹمی !!{formula}s نیں تے ایہ !!{operator}s تحت بند اے۔
  $\Frm[L_0]$ ایہو جہی سب توں چھوٹی جماعت اے، سو
  $\Frm[L_0] \subseteq S$۔ ایس لئی $\Frm[L_0] = S$، تے ہر فارمولے
  وچ خاصیت~$P$ اے۔
\end{proof}

\begin{prop}\ollabel{prop:balanced}
   $\Frm[L_0]$ وچ ہر !!{formula} \emph{متوازن} اے، یعنی اوہدے وچ
   کھبیاں قوساں اونّیاں ای نیں جِنّیاں سجیاں قوساں۔
\end{prop}

\begin{prob}
\olref[pl][syn][pre]{prop:balanced} ثابت کرو۔
\end{prob}

\begin{prop} \ollabel{prop:noinit}
کسی !!a{formula} دا پورے توں چھوٹا مُڈھلا ٹکڑا !!a{formula} نہیں ہُندا۔
\end{prop}

\begin{prob}
  \olref[pl][syn][pre]{prop:noinit} ثابت کرو۔
\end{prob}

\begin{prop}[اکّو طرح نال پڑھت]
ہر !!{formula}~$!A$، جیہڑا $ \Frm[L_0]$ وچ اے، اوہدا اگلیاں صورتاں
وچوں کسے اک دے طور تے تجزیہ ٹھیک اکّو طریقے نال ہُندا اے:
\begin{enumerate}
\tagitem{prvFalse}{$\lfalse$۔}{}

\tagitem{prvTrue}{$\ltrue$۔}{}

\item $\Obj p_n$، کسے $\Obj p_n \in  \PVar$ لئی۔

\tagitem{prvNot}{$\lnot !B$، کسے !!{formula}~$!B$ لئی۔}{}

\tagitem{prvAnd}{$(!B \land !C)$، کسے !!{formula}s $!B$ تے~$!C$ لئی۔}{}

\tagitem{prvOr}{$(!B \lor !C)$، کسے !!{formula}s $!B$ تے~$!C$ لئی۔}{}

\tagitem{prvIf}{$(!B \lif !C)$، کسے !!{formula}s $!B$ تے~$!C$ لئی۔}{}

\tagitem{prvIff}{$(!B \liff !C)$، کسے !!{formula}s $!B$ تے~$!C$ لئی۔}{}
\end{enumerate}
ہور ایہ تجزیاتی پڑھت \emph{اکّو ای} ہُندی اے۔
\end{prop}

\begin{proof}
$!A$ اُتے استقرا نال۔ مثال لئی، منّو کہ $!A$ دیاں دو وکھریاں پڑھتاں
$(!B \lif !C)$ تے $(!B' \lif !C')$ نیں۔ تاں $!B$ تے $!B'$ لازماً
اکّو ہون گے (نہیں تاں اک دوجے دا پورے توں چھوٹا مُڈھلا ٹکڑا ہُندا)؛ سو
جے $!A$ دیاں دونیں پڑھتاں وکھریاں نیں تاں وجہ ایہو ہو سکدی اے کہ
$!C$ تے $!C'$ علامتاں دی اکّو لڑی دیاں وکھریاں پڑھتاں نیں، جیہڑی گل
استقرائی مفروضے موجب ناممکن اے۔
\end{proof}

\begin{defn}[یکساں بدل]
جے $!A$ تے $!B$، !!{formula}s نیں، تے $\Obj p_i$ اک !!{propositional
variable} اے، تاں $\Subst{!A}{!B}{\Obj p_i}$ اوہ نتیجہ ظاہر کردا اے
جیہڑا $\Obj p_i$ دی ہر واری آمد نوں $!B$ دی اک آمد نال $!A$ وچ بدل
کے ملدا اے؛ ایسّے طرح $\Obj p_1$، \dots،~$\Obj p_n$ نوں اکّو ویلے
!!{formula}s $!B_1$، \dots،~$!B_n$ نال بدلنا
$\SSubst{!A}{\subst{!B_1}{\Obj p_1},\dots,\subst{!B_n}{\Obj p_n}}$
نال ظاہر کیتا جاندا اے۔
\end{defn}

\begin{prob} تھلے دتے پنجاں !!{formula}s وچوں ہر اک لئی طے کرو کہ کی
 اوہ !!{formula} بدل \( \Subst{!A}{!B}{\Obj p_i} \) دے روپ وچ لکھیا
 جا سکدا اے، جتھے \( !A \) ایہ اے: (i) \( \Obj p_0 \)؛ (ii)
 \( ( \lnot \Obj p_0 \land \Obj p_1) \)؛ تے (iii)
 \( ( ( \lnot \Obj p_0 \lif \Obj p_1 ) \land \Obj p_2 ) \)۔ ہر
 صورت وچ متعلق بدل صاف صاف دسو۔
  \begin{enumerate}
    \item \( \Obj p_1 \)
    \item \( ( \lnot \Obj p_0 \land \Obj p_0 ) \)
    \item \( ( ( \Obj p_0 \lor \Obj p_1 ) \land \Obj p_2 )  \)
    \item \( \lnot ( ( \Obj p_0 \lif \Obj p_1 ) \land \Obj p_2 ) \)
    \item \( (( \lnot ( \Obj p_0 \lif \Obj p_1 ) \lif ( \Obj p_0 \lor \Obj p_1 )) \land \lnot ( \Obj p_0 \land \Obj p_1 )) \)
  \end{enumerate}
\end{prob}

\begin{prob}
استقرا نال $\Subst{!A}{!B}{p}$ دی ریاضیاتی طور تے ٹھوس تعریف دیو۔
\end{prob}

\end{document}
